\documentclass[12 pt]{extarticle}

	\usepackage[french]{babel}
	\usepackage[utf8]{inputenc}  
	\usepackage[T1]{fontenc}
	\usepackage{amssymb}
	\usepackage[mathscr]{euscript}
	\usepackage{stmaryrd}
	\usepackage{amsmath}
	\usepackage{tikz}
	\usepackage[all,cmtip]{xy}
	\usepackage{amsthm}
	\usepackage{varioref}
	\usepackage{geometry}
	\geometry{a4paper}
	\usepackage{lmodern}
	\usepackage{hyperref}
	\usepackage{array}
	 \usepackage{fancyhdr}
\renewcommand{\theenumi}{\alph{enumi})}
	\pagestyle{empty}
	\theoremstyle{plain} 
	\geometry{
 a4paper,
 total={170mm,267mm},
 left=15mm,
 top=10mm,
 bottom=10mm,
 right=12mm
 }
	\usepackage{multicol}
	
	\title{Exercices Chapitre 1}
	\date{}
	\begin{document}

\begin{center}{\Large Chapitre 8- Angles}\\ 
 \end{center}
 
 \subsection*{Exercice 0 - Tracé de triangles}
 On rappelle que la somme des angles d'un triangle est $180^o$.
Tracer des triangles avec les conditions suivantes :
\begin{multicols}{2} \begin{enumerate}
\item $AB= 3~cm$, $AC=4~cm$ et $BC = 5~cm$.
\item $AB= 5~cm$, $\widehat{BAC} = 30^o$ et $\widehat{ABC}= 60^o$.
\item $AB=BC = 4~cm$ et $\widehat{ABC}= 120^o$. 
\item $AB= 3~cm$, $AC=6~cm$ et $\widehat{BAC}= 70^o$.
\item $AB = 4~cm$, $\widehat{BAC} = 80^o$ et $\widehat{ACB}= 70^o$.
\item $\widehat{ABC}= 75^o$, $AB= 6~cm$ et $AC= 7~cm$.
\item $AB = 7~cm$, $BC=6~cm$ et $\widehat{BAC}= 45^o$.
\item $AB= 3~cm$, $AC=8~cm$ et $BC = 4~cm$.

\end{enumerate} \end{multicols}
 
 \subsection*{Exercice 1}
 Deux angles $\widehat{AOB}$ et $\widehat{BOC}$ sont adjacents et $[OM)$ est la bissectrice de l'angle 
$\widehat{BOC}$. \begin{enumerate}
\item Construire la figure sachant que $\widehat{AOB}=48^o$ et $\widehat{AOC}= 112^o$. Calculer la mesure de
l'angle $\widehat{AOM}$. 
\item Si $\alpha$ et $\beta$ désignent les mesures des angles $\widehat{AOB}$ et $\widehat{AOC}$, montrer que l'on a toujours $\widehat{AOM}= \frac12(\alpha+\beta)$.  
\end{enumerate}
  
\subsection*{Exercice 2 }
  Soient $[OM)$ et $[ON)$ les bissectrices de deux angles adjacents $\widehat{AOB}$ et $\widehat{AOC}$. 
\begin{enumerate}
\item Construire la figure pour $\widehat{AOB}= 64^o$ et $\widehat{AOC}= 108^o$.
\item Calculer la mesure de l'angle $\widehat{MON}$. Comparer le résultat à la mesure de $\widehat{BOC}$. 
\item Énoncer un théorème donnant la mesure de l'angle des bissectrices de deux angles adjacents de mesures $\alpha$ et $\beta$.
\end{enumerate}
\subsection*{Exercice 3}
On construit un angle $\widehat{AOB}= 72^o$, puis les angles droits $\widehat{AOA'}$ et $\widehat{BOB'}$ adjacents au premier. 
\begin{enumerate}
\item Calculer la mesure de l'angle $\widehat{A'OB'}$. Comment sont les deux angles $\widehat{AOB}$ et $\widehat{A'OB'}$ ? 
\item Montrer que les bissectrices $[OM)$ et $[ON)$ des angles $\widehat{AOB}$ et $\widehat{A'OB'}$ sont alignées.
\item Calculer l'angle des bissectrices des angles droits $\widehat{AOA'}$ et $\widehat{BOB'}$.
\end{enumerate}

\subsection*{Exercice 4}
Soit un angle $\widehat{AOB}$ de $54^o$. On construit les angles droits $\widehat{AOA'}$ et $\widehat{BOB'}$ non adjacents à l'angle $\widehat{AOB}$.\begin{enumerate}
\item Montrer que les angles $\widehat{AOB}$ et $\widehat{A'OB'}$ sont supplémentaires et qu'ils ont même bissectrice $[OM)$.
\item Démontrer que les angles $\widehat{AOB'}$ et $\widehat{A'OB}$ sont égaux. Calculer l'angle de leurs bissectrices.
\end{enumerate}

\subsection*{Exercice 5}

On construit trois angles successivement adjacents 
$\widehat{AOC} = 48^o$, $\widehat{COD}= 72^o$ et $\widehat{DOB}= 34^o$. Puis on trace les demi-droites $[OM)$, $[ON)$, $[OP)$ et $[OQ)$ bissectrices des angles $\widehat{AOC}$, $\widehat{AOD}$, $\widehat{BOC}$ et $\widehat{BOD}$. \begin{enumerate}
\item Calculer les angles $\widehat{MON}$ et $\widehat{POQ}$ et comparer leur valeur à celle de l'angle $\widehat{COD}$.
\item Montrer que les angles $\widehat{MOQ}$ et $\widehat{NOP}$ ont même bissectrice. 
\end{enumerate}



 	\end{document}
